\title{DreamBooth: Fine Tuning Text-to-Image Diffusion Models for Subject-Driven Generation}

\begin{abstract}
Large text-to-image models achieved a remarkable leap in the evolution of AI, enabling high-quality and diverse synthesis of images from a given text prompt. However, these models lack the ability to mimic the appearance of subjects in a given reference set and synthesize novel renditions of them in different contexts. In this work, we present a new approach for ``personalization'' of text-to-image diffusion models. Given as input just a few images of a subject, we fine-tune a pretrained text-to-image model such that it learns to bind a unique identifier with that specific subject. Once the subject is embedded in the output domain of the model, the unique identifier can be used to synthesize novel photorealistic images of the subject contextualized in different scenes. By leveraging the semantic prior embedded in the model with a new autogenous class-specific prior preservation loss, our technique enables synthesizing the subject in diverse scenes, poses, views and lighting conditions that do not appear in the reference images. We apply our technique to several previously-unassailable tasks, including subject recontextualization, text-guided view synthesis, and artistic rendering, all while preserving the subject's key features. We also provide a new dataset and evaluation protocol for this new task of subject-driven generation. Project page:~\small{\url{https://dreambooth.github.io/}}
\end{abstract}

\section{Introduction}

Can you imagine your own dog traveling around the world, or your favorite bag displayed in the most exclusive showroom in Paris? What about your parrot being the main character of an illustrated storybook? Rendering such imaginary scenes is a challenging task that requires synthesizing instances of specific subjects (e.g., objects, animals) in new contexts such that they naturally and seamlessly blend into the scene.

Recently developed large text-to-image models have shown unprecedented capabilities, by enabling high-quality and diverse synthesis of images based on a text prompt written in natural language \cite{saharia2022photorealistic,ramesh2022hierarchical}. One of the main advantages of such models is the strong semantic prior learned from a large collection of image-caption pairs. Such a prior learns, for instance, to bind the word ``dog" with various instances of dogs that can appear in different poses and contexts in an image. While the synthesis capabilities of these models are unprecedented, they lack the ability to mimic the appearance of subjects in a given reference set, and synthesize novel renditions of the \emph{same subjects} in different contexts. The main reason is that the expressiveness of their output domain is limited;  even the most detailed textual description of an object may yield instances with different appearances. Furthermore, even models whose text embedding lies in a shared language-vision space \cite{radford2021learning} cannot accurately reconstruct the appearance of given subjects but only create variations of the image content (Figure~\ref{fig:explanatory_figure}). 

In this work, we present a new approach for ``personalization'' of text-to-image diffusion models (adapting them to user-specific image generation needs). Our goal is to expand the language-vision dictionary of the model such that it binds new words with specific subjects the user wants to generate. Once the new dictionary is embedded in the model, it can use these words to synthesize novel photorealistic images of the subject, contextualized in different scenes, while preserving their key identifying features. The effect is akin to a ``magic photo booth''---once a few images of the subject are taken, the booth generates photos of the subject in different conditions and scenes, as guided by simple and intuitive text prompts (Figure~\ref{fig:teaser}).

More formally, given a few images of a subject ($\sim$3-5), our objective is to implant the subject into the output domain of the model such that it can be synthesized with a \emph{unique identifier}. To that end, we propose a technique to represent a given subject with rare token identifiers and fine-tune a pre-trained, diffusion-based text-to-image framework. 

We fine-tune the text-to-image model with the input images and text prompts containing a unique identifier followed by the class name of the subject (e.g., ``A [V] dog”). The latter enables the model to use its prior knowledge on the subject class while the class-specific instance is bound with the unique identifier. In order to prevent \textit{language drift}~\cite{Lee2019CounteringLD,lu2020countering} that causes the model to associate the class name (e.g., ``dog") with the specific instance, we propose an \textit{autogenous, class-specific prior preservation loss}, which leverages the semantic prior on the class that is embedded in the model, and encourages it to generate diverse instances of the same class as our subject.

We apply our approach to a myriad of text-based image generation applications including recontextualization of subjects, modification of their properties, original art renditions, and more, paving the way to a new stream of previously unassailable tasks. We highlight the contribution of each component in our method via ablation studies, and compare with alternative baselines and related work. We also conduct a user study to evaluate subject and prompt fidelity in our synthesized images, compared to alternative approaches.

To the best of our knowledge, ours is the first technique that tackles this new challenging problem of subject-driven generation, allowing users, from just a few casually captured images of a subject, synthesize novel renditions of the subject in different contexts while maintaining its distinctive features.

To evaluate this new task, we also construct a new dataset that contains various subjects captured in different contexts, and propose a new evaluation protocol that measures the subject fidelity and prompt fidelity of the generated results. We make our dataset and evaluation protocol publicly available on the project webpage.

\section{Related work}

\textbf{Image Composition. } Image composition techniques \cite{wu2019gp,cong2020dovenet,lin2018st} aim to clone a given subject into a new background such that the subject melds into the scene. To consider composition in novel poses, one may apply 3D reconstruction techniques \cite{mildenhall2020nerf,Barron_2021_ICCV,boss2022samurai,verbin2021refnerf,poole2022dreamfusion} which usually works on rigid objects and require a larger number of views. Some drawbacks include scene integration (lighting, shadows, contact) and the inability to generate novel scenes. In contrast, our approach enable generation of subjects in novel poses and new contexts.

\textbf{Text-to-Image Editing and Synthesis. } Text-driven image manipulation has recently achieved significant progress using GANs ~\cite{goodfellow2014generative,brock2018large,karras2021alias,karras2019style,karras2020analyzing} combined with image-text representations such as CLIP~\cite{radford2021learning}, yielding realistic manipulations using text~\cite{patashnik2021styleclip,gal2021stylegan, xia2021tedigan, abdal2021clip2stylegan,bau2021paint,mokady2022self}. These methods work well on structured scenarios (e.g. human face editing) and can struggle over diverse datasets where subjects are varied. Crowson et al. \cite{crowson2022vqgan} use VQ-GAN \cite{esser2021taming} and train over more diverse data to alleviate this concern. Other works~\cite{avrahami2022blended, kim2022diffusionclip} exploit the recent diffusion models \cite{ho2020denoising,sohl2015deep,song2019generative,ho2020denoising,song2020denoising,rombach2021highresolution,nichol2021improved,song2020improved,palette,saharia2021image}, which achieve state-of-the-art generation quality over highly diverse datasets, often surpassing GANs~\cite{dhariwal2021diffusion}. While most works that require only text are limited to global editing \cite{crowson2022vqgan, kwon2021clipstyler}, Bar-Tal et al.~\cite{bar2022text2live} proposed a text-based localized editing technique without using masks, showing impressive results. While most of these editing approaches allow modification of global properties or local editing of a given image, none enables generating novel renditions of a given subject in new contexts.

There also exists work on text-to-image synthesis \cite{ding2021cogview,hinz2020semantic,tao2020df, li2019controllable,li2019object,qiao2019learn,qiao2019mirrorgan,ramesh2021zero,zhang2018photographic,crowson2022vqgan, gafni2022make, rombach2021highresolution,jain2021dreamfields}. Recent large text-to-image models such as Imagen~\cite{saharia2022photorealistic}, DALL-E2~\cite{ramesh2022hierarchical}, Parti~\cite{yu2022scaling}, CogView2~\cite{ding2022cogview2} and Stable Diffusion~\cite{rombach2021highresolution} demonstrated unprecedented semantic generation. These models do not provide fine-grained control over a generated image and use text guidance only. Specifically, it is challenging or impossible to preserve the identity of a subject consistently across synthesized images.

\textbf{Controllable Generative Models. } There are various approaches to control generative models, where some of them might prove to be viable directions for subject-driven prompt-guided image synthesis. Liu et al.~\cite{liu2019more} propose a diffusion-based technique allowing for image variations guided by reference image or text. To overcome subject modification, several works~\cite{nichol2021glide, avrahami2022blendedlatent} assume a user-provided mask to restrict the modified area. Inversion~\cite{choi2021ilvr,dhariwal2021diffusion,ramesh2022hierarchical} can be used to preserve a subject while modifying context. Prompt-to-prompt~\cite{hertz2022prompt} allows for local and global editing without an input mask. These methods fall short of identity-preserving novel sample generation of a subject.

In the context of GANs, Pivotal Tuning~\cite{roich2021pivotal} allows for real image editing by finetuning the model with an inverted latent code anchor, and Nitzan et al.~\cite{nitzan2022mystyle} extended this work to GAN finetuning on faces to train a personalized prior, which requires around 100 images and are limited to the face domain. Casanova et al.~\cite{casanova2021instance} propose an instance conditioned GAN that can generate variations of an instance, although it can struggle with unique subjects and does not preserve all subject details. 

Finally, the concurrent work of Gal \emph{et al}\onedot\cite{gal2022image} proposes a method to represent visual concepts, like an object or a style, through new tokens in the embedding space of a frozen text-to-image model, resulting in small personalized token embeddings. While this method is limited by the expressiveness of the frozen diffusion model, our fine-tuning approach enables us to embed the subject within the model's output domain, resulting in the generation of novel images of the subject which preserve its key visual features.

\section{Method}
\label{sec:method}

Given only a few (typically 3-5) casually captured images of a specific subject, without any textual description, our objective is to generate new images of the subject with high detail fidelity and with variations guided by text prompts. Example variations include changing the subject location, changing subject properties such as color or shape, modifying the subject's pose, viewpoint, and other semantic modifications. We do not impose any restrictions on input image capture settings and the subject image can have varying contexts. We next provide some background on text-to-image diffusion models (Sec.~\ref{sec:t2i}), then present our fine-tuning technique to bind a unique identifier with a subject described in a few images (Sec.~\ref{sec:finetune}), and finally propose a class-specific prior-preservation loss that enables us to overcome language drift in our fine-tuned model (Sec.~\ref{sec:method_prior_pres}).

\subsection{Text-to-Image Diffusion Models}
\label{sec:t2i}
Diffusion models are probabilistic generative models that are trained to learn a data distribution by the gradual denoising of a variable sampled from a Gaussian distribution. Specifically, we are interested in a pre-trained text-to-image diffusion model $\hat\mathbf{x}_\theta$ that, given an initial noise map ${\boldsymbol{\epsilon}} \sim \mathcal{N}(\mathbf{0}, \mathbf{I})$ and a conditioning vector $\mathbf{c}=\Gamma(\mathbf{P})$ generated using a text encoder $\Gamma$ and a text prompt $\mathbf{P}$, generates an image $\mathbf{x}_{\text{gen}}=\hat\mathbf{x}_\theta({\boldsymbol{\epsilon}}, \mathbf{c})$. They are trained using a squared error loss to denoise a variably-noised image or latent code $\mathbf{z}_t \coloneqq \alpha_t \mathbf{x} + \sigma_t {\boldsymbol{\epsilon}}$ as follows:

\begin{equation}
\label{eq:diffusion}
    \mathbb{E}_{\mathbf{x},\mathbf{c},{\boldsymbol{\epsilon}},t}\!\left[w_t \|\hat\mathbf{x}_\theta(\alpha_t \mathbf{x} + \sigma_t {\boldsymbol{\epsilon}}, \mathbf{c}) - \mathbf{x} \|^2_2\right]
\end{equation}

where $\mathbf{x}$ is the ground-truth image, $\mathbf{c}$ is a conditioning vector (e.g., obtained from a text prompt), and $\alpha_t, \sigma_t, w_t$ are terms that control the noise schedule and sample quality, and are functions of the diffusion process time $t \sim \mathcal{U}([0, 1])$. A more detailed description is given in the supplementary material.  

\subsection{Personalization of Text-to-Image Models}
\label{sec:finetune}

Our first task is to implant the subject instance into the output domain of the model such that we can query the model for varied novel images of the subject. One natural idea is to fine-tune the model using the few-shot dataset of the subject. Careful care had to be taken when fine-tuning generative models such as GANs in a few-shot scenario as it can cause overfitting and mode-collapse - as well as not capturing the target distribution sufficiently well. There has been research on techniques to avoid these pitfalls~\cite{Robb2020FewShotAO,ojha2021few,li2020few,mo2020freeze,Wang_2020_CVPR}, although, in contrast to our work, this line of work primarily seeks to generate images that resemble the target distribution but has no requirement of subject preservation. With regards to these pitfalls, we observe the peculiar finding that, given a careful fine-tuning setup using the diffusion loss from Eq~\ref{eq:diffusion}, large text-to-image diffusion models seem to excel at integrating new information into their domain without forgetting the prior or overfitting to a small set of training images. 

\paragraph{Designing Prompts for Few-Shot Personalization} Our goal is to ``implant" a new (\textit{unique identifier}, subject) pair into the diffusion model's ``dictionary'' . In order to bypass the overhead of writing detailed image descriptions for a given image set we opt for a simpler approach and label all input images of the subject ``a [identifier] [class noun]", where [identifier] is a unique identifier linked to the subject and [class noun] is a coarse class descriptor of the subject (e.g. cat, dog, watch, etc.). The class descriptor can be provided by the user or obtained using a classifier. We use a class descriptor in the sentence in order to tether the prior of the class to our unique subject and find that using a wrong class descriptor, or no class descriptor increases training time and language drift while decreasing performance. In essence, we seek to leverage the model's prior of the specific class and entangle it with the embedding of our subject's unique identifier so we can leverage the visual prior to generate new poses and articulations of the subject in different contexts.

\paragraph{Rare-token Identifiers} We generally find existing English words (e.g. ``unique'', ``special'') suboptimal since the model has to learn to disentangle them from their original meaning and to re-entangle them to reference our subject. This motivates the need for an identifier that has a weak prior in both the language model and the diffusion model. A hazardous way of doing this is to select random characters in the English language and concatenate them to generate a rare identifier (e.g. ``xxy5syt00"). In reality, the tokenizer might tokenize each letter separately, and the prior for the diffusion model is strong for these letters. We often find that these tokens incur the similar weaknesses as using common English words. Our approach is to find rare tokens in the vocabulary, and then invert these tokens into text space, in order to minimize the probability of the identifier having a strong prior. We perform a rare-token lookup in the vocabulary and obtain a sequence of rare token identifiers $f(\hat{\mathbf{V}})$, where $f$ is a tokenizer; a function that maps character sequences to tokens and $\hat{\mathbf{V}}$ is the decoded text stemming from the tokens $f(\hat{\mathbf{V}})$. The sequence can be of variable length $k$, and find that relatively short sequences of $k=\{1, ..., 3\}$ work well. Then, by inverting the vocabulary using the de-tokenizer on $f(\hat{\mathbf{V}})$ we obtain a sequence of characters that define our unique identifier $\hat{\mathbf{V}}$. For Imagen, we find that using uniform random sampling of tokens that correspond to 3 or fewer Unicode characters (without spaces) and using tokens in the T5-XXL tokenizer range of $\{5000, ..., 10000\}$ works well.

\subsection{Class-specific Prior Preservation Loss}
\label{sec:method_prior_pres}
In our experience, the best results for maximum subject fidelity are achieved by fine-tuning all layers of the model. This includes fine-tuning layers that are conditioned on the text embeddings, which gives rise to the problem of \textit{language drift}. Language drift has been an observed problem in language models~\cite{Lee2019CounteringLD,lu2020countering}, where a model that is pre-trained on a large text corpus and later fine-tuned for a specific task progressively loses syntactic and semantic knowledge of the language. To the best of our knowledge, we are the first to find a similar phenomenon affecting diffusion models, where to model slowly forgets how to generate subjects of the same class as the target subject.

Another problem is the possibility of \textit{reduced output diversity}. Text-to-image diffusion models naturally posses high amounts of output diversity. When fine-tuning on a small set of images we would like to be able to generate the subject in novel viewpoints, poses and articulations. Yet, there is a risk of reducing the amount of variability in the output poses and views of the subject (e.g. snapping to the few-shot views). We observe that this is often the case, especially when the model is trained for too long.

To mitigate the two aforementioned issues, we propose an autogenous class-specific prior preservation loss that encourages diversity and counters language drift. In essence, our method is to supervise the model with its \textit{own generated samples}, in order for it to retain the prior once the few-shot fine-tuning begins. This allows it to generate diverse images of the class prior, as well as retain knowledge about the class prior that it can use in conjunction with knowledge about the subject instance. Specifically, we generate data $\mathbf{x}_\text{pr}=\hat{\mathbf{x}}(\mathbf{z}_{t_1}, \mathbf{c}_\text{pr})$ by using the ancestral sampler on the frozen pre-trained diffusion model with random initial noise $\mathbf{z}_{t_1} \sim \mathcal{N}(\mathbf{0}, \mathbf{I})$ and conditioning vector $\mathbf{c}_\text{pr} \coloneqq \Gamma(f(\text{"a [class noun]"}))$. The loss becomes:

\begin{multline}
\label{eq:prior_pres}
    \mathbb{E}_{\mathbf{x},\mathbf{c},{\boldsymbol{\epsilon}},{\boldsymbol{\epsilon}}^\prime,t}[w_t \|\hat\mathbf{x}_\theta(\alpha_t \mathbf{x} + \sigma_t {\boldsymbol{\epsilon}}, \mathbf{c}) - \mathbf{x} \|^2_2 + \\
    \lambda w_{t^\prime} \|\hat\mathbf{x}_\theta(\alpha_{t^\prime} \mathbf{x}_\text{pr} + \sigma_{t^\prime} {\boldsymbol{\epsilon}}^\prime, \mathbf{c}_\text{pr}) - \mathbf{x}_\text{pr} \|^2_2],
\end{multline}

where the second term is the prior-preservation term that supervises the model with its own generated images, and $\lambda$ controls for the relative weight of this term. Figure~\ref{fig:system_overview} illustrates the model fine-tuning with the class-generated samples and prior-preservation loss. Despite being simple, we find this prior-preservation loss is effective in encouraging output diversity and  in overcoming language-drift. We also find that we can train the model for more iterations without risking overfitting. We find that \(\sim \) 1000 iterations with $\lambda = 1$ and learning rate $10^{-5}$ for Imagen~\cite{saharia2022photorealistic} and $5 \times 10^{-6}$ for Stable Diffusion~\cite{rombach2022high}, and with a subject dataset size of 3-5 images is enough to achieve good results. During this process, $\sim 1000$ ``a [class noun]'' samples are generated - but less can be used. The training process takes about 5 minutes on one TPUv4 for Imagen, and 5 minutes on a NVIDIA A100 for Stable Diffusion.

\section{Experiments} 
\label{sec:app}

In this section, we show experiments and applications. Our method enables a large expanse of text-guided semantic modifications of our subject instances, including recontextualization, modification of subject properties such as material and species, art rendition, and viewpoint modification. Importantly, across all of these modifications, we are able to \textbf{preserve the unique visual features that give the subject its identity and essence}. If the task is recontextualization, then the subject features are unmodified, but appearance (e.g., pose) may change. If the task is a stronger semantic modification, such as crossing between our subject and another species/object, then the key features of the subject are preserved after modification. In this section, we reference the subject's unique identifier using [V]. We include specific Imagen and Stable Diffusion implementation details in the supp. material.

\subsection{Dataset and Evaluation}

\paragraph{Dataset}
We collected a dataset of 30 subjects, including unique objects and pets such as backpacks, stuffed animals, dogs, cats, sunglasses, cartoons, etc. We separate each subject into two categories: objects and live subjects/pets. 21 of the 30 subjects are objects, and 9 are live subjects/pets. We provide one sample image for each of the subjects in Figure~\ref{fig:dataset}. Images for this dataset were collected by the authors or sourced from Unsplash~\cite{unsplash}. We also collected 25 prompts: 20 recontextualization prompts and 5 property modification prompts for objects; 10 recontextualization, 10 accessorization, and 5 property modification prompts for live subjects/pets. The full list of prompts can be found in the supplementary material.

For the evaluation suite we generate four images per subject and per prompt, totaling 3,000 images. This allows us to robustly measure performances and generalization capabilities of a method. We make our dataset and evaluation protocol publicly available on the project webpage for future use in evaluating subject-driven generation.

\paragraph{Evaluation Metrics} One important aspect to evaluate is subject fidelity: the preservation of subject details in generated images. For this, we compute two metrics: CLIP-I and DINO~\cite{caron2021emerging}. CLIP-I is the average pairwise cosine similarity between CLIP~\cite{radford2021learning} embeddings of generated and real images. Although this metric has been used in other work~\cite{gal2022image}, it is not constructed to distinguish between different subjects that could have highly similar text descriptions (e.g. two different yellow clocks). Our proposed DINO metric is the average pairwise cosine similarity between the ViT-S/16 DINO embeddings of generated and real images. This is our preferred metric, since, by construction and in contrast to supervised networks, DINO is not trained to ignore differences between subjects of the same class. Instead, the self-supervised training objective encourages distinction of unique features of a subject or image. The second important aspect to evaluate is prompt fidelity, measured as the average cosine similarity between prompt and image CLIP embeddings. We denote this as CLIP-T.

\subsection{Comparisons}

We compare our results with Textual Inversion, the recent concurrent work of Gal et al.~\cite{gal2022image}, using the hyperparameters provided in their work. We find that this work is the only comparable work in the literature that is subject-driven, text-guided and generates novel images. We generate images for DreamBooth using Imagen, DreamBooth using Stable Diffusion and Textual Inversion using Stable Diffusion. We compute DINO and CLIP-I subject fidelity metrics and the CLIP-T prompt fidelity metric. In Table~\ref{table:comparison_experiment} we show sizeable gaps in both subject and prompt fidelity metrics for DreamBooth over Textual Inversion. We find that DreamBooth (Imagen) achieves higher scores for both subject and prompt fidelity than DreamBooth (Stable Diffusion), approaching the upper-bound of subject fidelity for real images. We believe that this is due to the larger expressive power and higher output quality of Imagen.

Further, we compare Textual Inversion (Stable Diffusion) and DreamBooth (Stable Diffusion) by conducting a user study. 

For subject fidelity, we asked 72 users to answer questionnaires of 25 comparative questions (3 users per questionnaire), totaling 1800 answers. Samples are randomly selected from a large pool. Each question shows the set of real images for a subject, and one generated image of that subject by each method (with a random prompt). Users are asked to answer the question: ``Which of the two images best reproduces the identity (e.g. item type and details) of the reference item?'', and we include a ``Cannot Determine / Both Equally'' option. Similarly for prompt fidelity, we ask ``Which of the two images is best described by the reference text?''. We average results using majority voting and present them in Table~\ref{table:user_study}. We find an overwhelming preference for DreamBooth for both subject fidelity and prompt fidelity. This shines a light on results in Table~\ref{table:comparison_experiment}, where DINO differences of around $0.1$ and CLIP-T differences of $0.05$ are significant in terms of user preference. Finally, we show qualitative comparisons in Figure~\ref{fig:comparison_gal}. We observe that DreamBooth better preserves subject identity, and is more faithful to prompts. We show samples of the user study in the supp. material.

\subsection{Ablation Studies}
\label{sec:app_ablation}

\paragraph{Prior Preservation Loss Ablation} We fine-tune Imagen on 15 subjects from our dataset, with and without our proposed prior preservation loss (PPL). The prior preservation loss seeks to combat language drift and preserve the prior. We compute a prior preservation metric (PRES) by computing the average pairwise DINO embeddings between generated images of random subjects of the prior class and real images of our specific subject. The higher this metric, the more similar random subjects of the class are to our specific subject, indicating collapse of the prior. We report results in Table~\ref{table:ablation_prior_pres} and observe that PPL substantially counteracts language drift and helps retain the ability to generate diverse images of the prior class. Additionally, we compute a diversity metric (DIV) using the average LPIPS~\cite{zhang2018perceptual} cosine similarity between generated images of same subject with same prompt. We observe that our model trained with PPL achieves higher diversity (with slightly diminished subject fidelity), which can also be observed qualitatively in Figure~\ref{fig:prior_pres_target}, where our model trained with PPL overfits less to the environment of the reference images and can generate the dog in more diverse poses and articulations.

\paragraph{Class-Prior Ablation} We finetune Imagen on a subset of our dataset subjects (5 subjects) with no class noun, a randomly sampled incorrect class noun, and the correct class noun. With the correct class noun for our subject, we are able to faithfully fit to the subject, take advantage of the class prior, allowing us to generate our subject in various contexts. When an incorrect class noun (e.g. ``can'' for a backpack) is used, we run into contention between our subject and and the class prior - sometimes obtaining cylindrical backpacks, or otherwise misshapen subjects. If we train with no class noun, the model does not leverage the class prior, has difficulty learning the subject and converging, and can generate erroneous samples. Subject fidelity results are shown in Table~\ref{table:class_name}, with substantially higher subject fidelity for our proposed approach.

\subsection{Applications}

\paragraph{Recontextualization}
We can generate novel images for a specific subject in different contexts (Figure~\ref{fig:recontext}) with descriptive prompts (``a [V] [class noun] [context description]''). Importantly, we are able to generate the subject in new poses and articulations, with previously unseen scene structure and realistic integration of the subject in the scene (e.g. contact, shadows, reflections). 

\paragraph{Art Renditions} Given a prompt ``a painting of a [V] [class noun] in the style of [famous painter]'' or ``a statue of a [V] [class noun] in the style of [famous sculptor]'' we are able to generate artistic renditions of our subject. Unlike style transfer, where the source structure is preserved and only the style is transferred, we are able to generate meaningful, novel variations depending on the artistic style, while preserving subject identity. E.g, as shown in Figure~\ref{fig:applications}, ``Michelangelo'', we generated a pose that is novel and not seen in the input images.

\paragraph{Novel View Synthesis} We are able to render the subject under novel viewpoints. In Figure~\ref{fig:applications}, we generate new images of the input cat (with consistent complex fur patterns) under new viewpoints. We highlight that the model has not seen this specific cat from behind, below, or above - yet it is able to extrapolate knowledge from the class prior to generate these novel views given only 4 frontal images of the subject.

\paragraph{Property Modification} We are able to modify subject properties. For example, we show crosses between a specific Chow Chow dog and different animal species in the bottom row of Figure~\ref{fig:applications}. We prompt the model with sentences of the following structure: ``a cross of a [V] dog and a [target species]''. In particular, we can see in this example that the identity of the dog is well preserved even when the species changes - the face of the dog has certain unique features that are well preserved and melded with the target species. Other property modifications are possible, such as material modification (e.g. ``a transparent [V] teapot'' in Figure~\ref{fig:recontext}). Some are harder than others and depend on the prior of the base generation model.

\subsection{Limitations}
We illustrate some failure models of our method in Figure~\ref{fig:limitations}. The first is related to not being able to accurately generate the prompted context. Possible reasons are a weak prior for these contexts, or difficulty in generating both the subject and specified concept together due to low probability of co-occurrence in the training set. The second is context-appearance entanglement, where the appearance of the subject changes due to the prompted context, exemplified in Figure~\ref{fig:limitations} with color changes of the backpack. Third, we also observe overfitting to the real images that happen when the prompt is similar to the original setting in which the subject was seen.

Other limitations are that some subjects are easier to learn than others (e.g. dogs and cats). Occasionally, with subjects that are rarer, the model is unable to support as many subject variations. Finally, there is also variability in the fidelity of the subject and some generated images might contain hallucinated subject features, depending on the strength of the model prior, and the complexity of the semantic modification.

\section{Conclusions}
We presented an approach for synthesizing novel renditions of a subject using a few images of the subject and the guidance of a text prompt. Our key idea is to embed a given subject instance in the output domain of a text-to-image diffusion model by binding the subject to a unique identifier. Remarkably - this fine-tuning process can work given only 3-5 subject images, making the technique particularly accessible. We demonstrated a variety of applications with animals and objects in generated photorealistic scenes, in most cases indistinguishable from real images.

\end{document}
